% ============================================================
% 英文论文模板（IEEE 双栏，paper-en）
% ------------------------------------------------------------
% 使用 IEEEtran 文档类，版式与 IEEE 会议/期刊论文一致。
% 编译引擎：pdflatex（LaTeXcellence 已按此配置）。
% 修改指南：
%   1. \title / \author 处填写论文题目与作者信息；
%   2. 摘要在 abstract 环境，关键词在 \begin{IEEEkeywords}；
%   3. 各章节正文直接替换示例文字；
%   4. 参考文献使用文中 thebibliography 环境，
%      若习惯 BibTeX，可删除该环境改用
%      \bibliographystyle{IEEEtran} + \bibliography{refs}。
% ============================================================
\documentclass[conference]{IEEEtran}

% ---- 常用宏包 ----
\usepackage{amsmath, amssymb, amsfonts}   % 数学公式
\usepackage{graphicx}                     % 插图
\usepackage{booktabs}                     % 三线表
\usepackage{multirow}                     % 表格合并单元格
\usepackage{xcolor}                       % 颜色
\usepackage[colorlinks=true, linkcolor=black, citecolor=blue, urlcolor=blue]{hyperref}

% 文中自定义命令示例：向量用粗斜体
\newcommand{\vect}[1]{\boldsymbol{#1}}

\begin{document}

\title{An Example Paper Title: Written with the IEEEtran Class}

% \author 中的 \IEEEauthorblockN 为作者名、\IEEEauthorblockA 为单位信息
\author{
  \IEEEauthorblockN{First Author}
  \IEEEauthorblockA{Department of Computer Science, Example University\\
  City, Country\\
  Email: first.author@example.edu}
  \and
  \IEEEauthorblockN{Second Author}
  \IEEEauthorblockA{School of Artificial Intelligence, Example Institute\\
  City, Country\\
  Email: second.author@example.ac.cn}
}

\maketitle

\begin{abstract}
  Replace this paragraph with the abstract of your paper.
  The abstract should state the problem, the proposed method,
  and the main results in 150--250 words. This template uses the
  IEEEtran conference class and compiles with pdfLaTeX.
\end{abstract}

\begin{IEEEkeywords}
  keyword one, keyword two, keyword three, LaTeX, IEEEtran
\end{IEEEkeywords}

\section{Introduction}
% 引言：说明研究背景、现有不足与本文贡献
This is the introduction. Explain the background, the limitations of
existing work, and summarize your contributions as an itemized list:

\begin{itemize}
  \item \textbf{Contribution one:} a brief statement of the first contribution.
  \item \textbf{Contribution two:} a brief statement of the second contribution.
  \item \textbf{Contribution three:} extensive experiments on benchmark datasets.
\end{itemize}

\section{Related Work}
% 相关工作：分组综述代表性方法，并说明与本文的区别
Summarize the most related approaches and clarify how your work differs
from them. Cite references with \verb|\cite|, e.g., deep models~\cite{lecun2015}.

\section{Method}
% 方法：给出总体框架、符号定义与关键公式
Describe your overall framework here. An example equation:
\begin{equation}
  \mathcal{L}(\vect{\theta}) =
  \frac{1}{N} \sum_{i=1}^{N}
  \ell \bigl( f(\vect{x}_i; \vect{\theta}),\, y_i \bigr)
  + \lambda \, \lVert \vect{\theta} \rVert_2^2
  \label{eq:loss}
\end{equation}
where $\ell(\cdot)$ is the loss function, $\lambda$ balances the
regularization term, and $N$ is the number of training samples.
Reference the equation as \eqref{eq:loss}.

\section{Experiments}
% 实验：数据集、评价指标、对比结果（表）与定性分析（图）
\subsection{Datasets and Metrics}
Briefly describe the datasets and evaluation metrics.

\subsection{Main Results}
Table~\ref{tab:results} reports the comparison with baselines.

% 三线表示例：\toprule / \midrule / \bottomrule
\begin{table}[!t]
  \centering
  \caption{Comparison with baseline methods on the benchmark.}
  \label{tab:results}
  \begin{tabular}{lccc}
    \toprule
    Method & Acc.(\%) & F1 (\%) & Params (M) \\
    \midrule
    Baseline A        & 85.2 & 84.1 & 21.3 \\
    Baseline B        & 87.6 & 86.9 & 34.8 \\
    Ours              & \textbf{90.4} & \textbf{89.7} & 28.5 \\
    \bottomrule
  \end{tabular}
\end{table}

% 插图示例：把 image.png 换成你的图片文件
% \begin{figure}[!t]
%   \centering
%   \includegraphics[width=\linewidth]{image.png}
%   \caption{An example figure caption.}
%   \label{fig:framework}
% \end{figure}

\section{Conclusion}
% 结论：总结工作并展望未来方向
Restate the problem and your approach, highlight the key findings,
and outline future work.

% ---- 参考文献：按正文引用顺序排列 ----
\begin{thebibliography}{9}
\bibitem{lecun2015}
Y. LeCun, Y. Bengio, and G. Hinton,
``Deep learning,''
\emph{Nature}, vol. 521, no. 7553, pp. 436--444, 2015.

\bibitem{vaswani2017}
A. Vaswani \emph{et al.},
``Attention is all you need,''
in \emph{Proc. Adv. Neural Inf. Process. Syst. (NeurIPS)}, 2017, pp. 5998--6008.
\end{thebibliography}

\end{document}
